\section{Monomorfizmy, epimorfizmy, izomorfizmy}

Przypomnijmy najpierw pojęcia izomorfizmu, mono- i epimorfizmu wprowadzone pod koniec poprzedniego rozdziału; zobaczymy je teraz w działaniu na przykładach.

\begin{definicja}
Niech $f\colon A\to B$ będzie strzałką w kategorii $\C$. Strzałkę $f$ nazywamy \emph{izomorfizmem}, jeśli istnieje $g\colon B\to A$ w $\C$ taka, że
\[
g\circ f = \id_A \quad\text{oraz}\quad f\circ g = \id_B.
\]
\end{definicja}

\begin{definicja}
Strzałkę $f\colon A\to B$ nazywamy \emph{monomorfizmem} (\emph{mono}), jeśli dla każdych $g,h\colon C\to A$ w $\C$
\[
f\circ g = f\circ h \;\Longrightarrow\; g=h,
\]
tzn.\ diagram
\[
\begin{tikzcd}
C \arrow[r, shift left, "g"] \arrow[r, shift right, "h"'] & A \arrow[r, "f"] & B
\end{tikzcd}
\]
daje implikację $f\circ g = f\circ h \Rightarrow g=h$.
\end{definicja}

\begin{definicja}
Strzałkę $f\colon A\to B$ nazywamy \emph{epimorfizmem} (\emph{epi}), jeśli dla każdych $i,j\colon B\to C$ w $\C$
\[
i\circ f = j\circ f \;\Longrightarrow\; i=j.
\]
\end{definicja}

\begin{stwierdzenie}
Funkcja $f\colon A\to B$ jest mono w $\Set$ wtedy i tylko wtedy, gdy jest różnowartościowa.
\end{stwierdzenie}

\begin{stwierdzenie}
$f\colon A\to B$ jest mono w $\C$ wtedy i tylko wtedy, gdy $f\colon B\to A$ jest epi w $\C^{\op}$.
\end{stwierdzenie}

\begin{proof}
Niech $f\colon A\to B$ będzie mono w $\C$. W szczególności oznacza to, że $f\colon B\to A$ jest strzałką w $\C^{\op}$. Weźmy $i,j\colon A\to C$ w $\C^{\op}$ i załóżmy, że w $\C^{\op}$ zachodzi
\[
i\circ_{\C^{\op}} f = j\circ_{\C^{\op}} f,
\]
gdzie $f\colon A\to B$, $i,j\colon C\to A$ w $\C$. Z definicji składania w $\C^{\op}$ wynika, że w $\C$ zachodzi
\[
f\circ i = f\circ j,
\]
gdzie $f\colon A\to B$ oraz $i,j\colon C\to A$ w $\C$. Ponieważ $f$ jest mono w $\C$, dostajemy $i=j$ w $\C$. Ale to znaczy, że $i=j$ w $\C^{\op}$. Stąd $f$ jest epi w $\C^{\op}$. Implikacja w przeciwną stronę jest analogiczna.
\end{proof}

\subsection{Przykłady}

\begin{przyklad}
W kategorii $\Set$ mono $=$ ,,różnowartościowe'', a epi $=$ ,,na''.
\end{przyklad}

\begin{przyklad}
Niech $\C$ będzie kategorią indukowaną przez poset $(P,\leq)$. Weźmy $p,q\in P$. Wtedy dla strzałek
\[
\begin{tikzcd}
g,h\colon r \arrow[r] & p \arrow[r, "f"] & q
\end{tikzcd}
\]
warunek $f\circ g = f\circ h$ daje $g=h$ (gdyż wszystkie strzałki między tą samą parą obiektów są równe).
\end{przyklad}

\begin{wniosek}
Każda strzałka $f\colon p\leq q$ jest mono (i epi) w kategorii indukowanej przez poset $(P,\leq)$.
\end{wniosek}

\begin{przyklad}
Rozważmy kategorię $\Mon$ -- kategorię monoidów z homomorfizmami monoidów jako strzałkami. W tej kategorii rozważmy strzałkę
\[
f\colon (\mathbb{N},+,0) \longrightarrow (\mathbb{Z},+,0), \qquad f(i) = i.
\]
Ta strzałka jest:
\begin{enumerate}
\item mono,
\item epimorfizmem.
\end{enumerate}

\emph{Dowód (1).} Strzałka $f$ jest różnowartościowa, więc jest mono w $\Set$ -- a zatem także w $\Mon$.

\emph{Dowód (2).} Musimy pokazać, że dla każdej pary
\[
i,j\colon (\mathbb{Z},+,0) \longrightarrow M
\]
(gdzie $M=(M,\cdot,1)$ jest monoidem) zachodzi
\[
i\circ f = j\circ f \;\Longrightarrow\; i=j.
\]
Załóżmy, że $i\circ f = j\circ f$. Z tej równości wynika w szczególności, że
\[
i|_{\mathbb{N}} = j|_{\mathbb{N}}.
\]
Wystarczy więc pokazać, że $i(-n) = j(-n)$ dla każdego $n\in\mathbb{N}$. Z faktu, że $i$ jest homomorfizmem mamy
\[
1 = i(0) = i(n + (-n)) = i(n)\cdot i(-n)
\qquad\text{oraz}\qquad
1 = i(0) = i((-n) + n) = i(-n)\cdot i(n),
\]
czyli $i(-n)$ jest obustronną odwrotnością elementu $i(n)$. W monoidzie element ma co najwyżej jedną obustronną odwrotność: jeśli $x\cdot m = 1 = m\cdot y$, to $x = x\cdot(m\cdot y) = (x\cdot m)\cdot y = y$. Analogicznie $j(-n)$ jest obustronną odwrotnością $j(n)$. Z $i|_{\mathbb{N}}=j|_{\mathbb{N}}$ wynika $i(n)=j(n)$, więc z jednoznaczności odwrotności $i(-n) = j(-n)$, a zatem $i=j$. \qed
\end{przyklad}

\begin{uwaga}
Powyższy przykład pokazuje, że strzałka będąca jednocześnie mono i epi nie musi być izomorfizmem ($f$ nie jest ,,na'', więc nie ma odwrotności). W $\Set$ tak być nie może: mono $+$ epi $=$ różnowartościowa i ,,na'' $=$ bijekcja $=$ izo. Inny przykład: w $\Top$ włożenie $\mathbb{Q} \hookrightarrow \mathbb{R}$ jest mono (jest różnowartościowe) i epi (dwa przekształcenia ciągłe o wartościach w przestrzeni Hausdorffa, zgodne na gęstym podzbiorze, są równe), ale nie jest homeomorfizmem.
\end{uwaga}

\section{Obiekty początkowe i końcowe}

W $\Set$:
\[
\begin{tikzcd}
\emptyset \arrow[r, "i"] & A \arrow[r] & \mathbf{1}
\end{tikzcd}
\]
$\emptyset$ jest obiektem początkowym, $\mathbf{1}$ (zbiór jednoelementowy) jest obiektem końcowym.

\begin{definicja}
Niech $\C$ będzie kategorią. Obiekt $0\in\C$ nazywamy \emph{początkowym}, jeśli dla każdego obiektu $A\in\C$ istnieje dokładnie jedna strzałka
\[
i_A\colon 0 \longrightarrow A.
\]
\end{definicja}

\begin{definicja}
Obiekt $\mathbf{1}\in\C$ nazywamy \emph{końcowym}, jeśli dla każdego $A\in\C$ istnieje dokładnie jedna strzałka
\[
!_A\colon A \longrightarrow \mathbf{1}.
\]
\end{definicja}

\begin{uwaga}
$0\in\C$ jest obiektem początkowym w $\C$ wtedy i tylko wtedy, gdy $0$ jest końcowym w $\C^{\op}$.
\end{uwaga}

\begin{stwierdzenie}
Obiekt początkowy (końcowy) jest wyznaczony jednoznacznie z dokładnością do izomorfizmu.
\end{stwierdzenie}

\begin{proof}
Załóżmy, że $0$, $0'$ są początkowe w kategorii $\C$. Wtedy istnieją jednoznacznie wyznaczone strzałki
\[
\begin{tikzcd}
0 \arrow[r, shift left, "u"] & 0' \arrow[l, shift left, "v"]
\end{tikzcd}
\]
spełniające:
\begin{enumerate}
\item $v\circ u = \id_0$ (gdyż obie strzałki $0\to 0$ muszą być równe, a $\id_0$ jest jedyną),
\item $u\circ v = \id_{0'}$ (analogicznie).
\end{enumerate}
Stąd $0\cong 0'$. Dowód dla obiektu końcowego jest analogiczny.
\end{proof}

\subsection{Przykłady}

\begin{przyklad}
W $\Set$: obiekt początkowy $=\emptyset$, obiekt końcowy $=$ dowolny zbiór jednoelementowy.
\end{przyklad}

\begin{przyklad}
W $\Grp$ grupa trywialna $\{e\}$ jest jednocześnie obiektem początkowym (jedyny homomorfizm $\{e\} \to G$ posyła $e$ na element neutralny) i końcowym (jedyny homomorfizm $G \to \{e\}$ jest stały). Obiekt o obu tych własnościach nazywamy \emph{obiektem zerowym}. W $\Set$ obiekty początkowy i końcowy są natomiast różne: $\emptyset \neq \{*\}$.
\end{przyklad}

\begin{przyklad}
Niech $\C$ będzie kategorią indukowaną przez poset $(P,\leq)$. Czym (jeśli istnieją) są obiekty początkowe i końcowe w $\C$?

\emph{Obiekt początkowy:} $0\leq p$ dla każdego $p\in P$. Widzimy, że $0$ jest elementem najmniejszym w $(P,\leq)$. Np.\ w
\[
\begin{tikzcd}
\bullet & & \bullet
\end{tikzcd}
\]
nie ma obiektu początkowego.

\emph{Obiekt końcowy:} $p \leq \mathbf{1}$ dla każdego $p\in P$, czyli $\mathbf{1}$ jest elementem największym w $(P,\leq)$.
\end{przyklad}

\begin{przyklad}
Rozważmy kategorię algebr typu $(1,0)$. Obiektami w tej kategorii są
\[
(A,s,a), \qquad \text{gdzie } s\colon A\to A,\; a\in A.
\]
Strzałkami w tej kategorii są
\[
f\colon (A,s,a) \longrightarrow (B,s',b),
\]
gdzie $f\colon A\to B$ jest funkcją spełniającą:
\begin{enumerate}
\item dla każdego $x\in A$: $s'(f(x)) = f(s(x))$,
\item $f(a)=b$.
\end{enumerate}
\end{przyklad}

\begin{stwierdzenie}
W tej kategorii obiektem początkowym jest algebra
\[
(\mathbb{N},\mathrm{succ},0), \qquad \mathrm{succ}(n)\stackrel{\mathrm{df}}{=} n+1.
\]
\end{stwierdzenie}

\begin{proof}[Szkic dowodu]
Niech $(B, s', b)$ będzie dowolną algebrą typu $(1,0)$. Strzałka $f\colon (\mathbb{N},\mathrm{succ},0)\to (B,s',b)$ musi spełniać
\[
f(0) = b, \qquad f(n+1) = f(\mathrm{succ}(n)) = s'(f(n)),
\]
a te dwa warunki \emph{definiują} $f$ rekurencyjnie -- strzałka więc istnieje, a jej jednoznaczność pokazujemy przez indukcję po $n$. Inicjalność $(\mathbb{N},\mathrm{succ},0)$ to zatem dokładnie zasada definiowania przez rekursję.
\end{proof}

\section{Produkty}

W $\Set$:
\[
A\times B = \{(a,b) : a\in A,\; b\in B\}.
\]
Zauważmy, że każdy produkt kartezjański $A\times B$ jest w parze z dwoma rzutowaniami
\[
\begin{tikzcd}
A & A\times B \arrow[l, "\pi_1"'] \arrow[r, "\pi_2"] & B
\end{tikzcd}
\]
gdzie $\pi_1(a,b) = a$, $\pi_2(a,b)=b$. Ponadto dla dowolnego zbioru $X$ wraz ze strzałkami $c_1\colon X\to A$, $c_2\colon X\to B$ istnieje jednoznacznie wyznaczona strzałka $X\to A\times B$ dana wzorem $x\mapsto (c_1(x),c_2(x))$.

\begin{definicja}
Niech $\C$ będzie kategorią. \emph{Diagramem produktu} dla obiektów $A,B\in\C$ nazywamy obiekt $P$ wraz ze strzałkami
\[
\begin{tikzcd}
A & P \arrow[l, "p_1"'] \arrow[r, "p_2"] & B
\end{tikzcd}
\]
takimi, że dla każdego obiektu $X\in\C$ wraz ze strzałkami
\[
\begin{tikzcd}
A & X \arrow[l, "x_1"'] \arrow[r, "x_2"] & B
\end{tikzcd}
\]
istnieje jednoznaczna strzałka $u\colon X\to P$ czyniąca następujący diagram przemiennym:
\[
\begin{tikzcd}
& X \arrow[ld, "x_1"'] \arrow[rd, "x_2"] \arrow[d, dashed, "\exists! u"] & \\
A & P \arrow[l, "p_1"] \arrow[r, "p_2"'] & B
\end{tikzcd}
\]
\end{definicja}

\begin{definicja}
Jeśli $P$ wraz z $A\xleftarrow{p_1} P\xrightarrow{p_2} B$ jest diagramem produktowym dla obiektów $A,B$, to $P$ nazywamy \emph{produktem} $A$ oraz $B$.
\end{definicja}

\begin{stwierdzenie}
Produkty są określone jednoznacznie z dokładnością do izomorfizmu.
\end{stwierdzenie}

\begin{proof}
Załóżmy, że
\[
\begin{tikzcd}
A & P \arrow[l, "p_1"'] \arrow[r, "p_2"] & B
\end{tikzcd}
\qquad\text{oraz}\qquad
\begin{tikzcd}
A & Q \arrow[l, "q_1"'] \arrow[r, "q_2"] & B
\end{tikzcd}
\]
są diagramami produktowymi dla $A,B$. $Q$ jest produktem, więc istnieje jednoznacznie określona strzałka $i\colon P\to Q$ taka, że
\[
q_1\circ i = p_1 \quad\text{oraz}\quad q_2\circ i = p_2.
\]
Analogicznie, ponieważ $P$ jest produktem, istnieje jednoznacznie określona strzałka $j\colon Q\to P$ taka, że
\[
q_1 = p_1\circ j, \qquad q_2 = p_2\circ j.
\]
Stąd $p_1\circ j\circ i = p_1$ oraz $p_2\circ j\circ i = p_2$. Z jednoznaczności strzałki czyniącej diagram produktowy przemiennym (a $\id_P$ także go uzupełnia) mamy $j\circ i = \id_P$. Analogicznie $i\circ j = \id_Q$.
\end{proof}
